\section{Agent architectures}
\label{sec:agent-architectures}
\slidesrc{Agent architectures: 1--41}

How is an agent built \emph{inside}? This chapter (slides from a course by URV, Tarragona,
with comments by prof.\ Mascardi) first looks at what an agent must cope with, namely its
\emph{environment}, and then at the three main families of \textbf{agent architectures}:
\begin{enumerate}[nosep]
  \item \textbf{reactive}: fast stimulus--response rules, no model of the world (Brooks);
  \item \textbf{deliberative}: an explicit symbolic model of the world and a planner;
  \item \textbf{hybrid}: a reactive layer and a deliberative layer together.
\end{enumerate}
The BDI architecture of Part~V is a practical-reasoning architecture that combines
deliberation (goals, plans) with reactivity (events), in the spirit of the hybrid approach.

\subsection{Agents and their environments}
\label{sec:agents-and-environments}
\slidesrc{Agent architectures: 1--10}

\subsubsection{Three classic definitions of agent}
\slidesrc{Agent architectures: 3--5}

The slides recall three classic definitions (the sources are added):
\begin{enumerate}
  \item ``An agent is anything that can be viewed as \emph{perceiving its environment
        through sensors} and \emph{acting upon that environment through effectors}.''
        \hfill (Russell and Norvig)
  \item ``Autonomous agents are computational systems that inhabit some complex dynamic
        environment, sense and act autonomously in this environment, and by doing so realize
        a set of goals or tasks for which they are designed.'' \hfill (P.~Maes)
  \item ``An autonomous agent is a system situated within and a part of an environment that
        senses that environment and acts on it, over time, in pursuit of its own agenda and
        so as to affect what it senses in the future.'' \hfill (Franklin and Graesser)
\end{enumerate}
All three share the \textbf{basic abstract view} of slide 4: a loop between agent and
environment, through sensors (input) and effectors (output); the same picture as
Section~\ref{sec:ai-definition}. The third definition adds an important nuance:
the agent's actions change what it will \emph{perceive in the future}.

\begin{nbdef}{Reactivity (slide 5)}
An agent has to be able to \textbf{react}, i.e.\ adapt its behaviour, in an appropriate way
to the dynamic changes in its environment: other computational agents, human
agents/users, external information sources (e.g.\ sensors), physical objects (e.g.\
robots), the Internet, \ldots{} Reactivity is one of several properties an intelligent agent
should have (Section~\ref{sec:what-is-an-agent}).
\end{nbdef}

\subsubsection{Kinds of environments}
\slidesrc{Agent architectures: 6--10}

The difficulty of building an agent depends heavily on the environment it lives in. Five
dimensions classify environments; in each pair, the first is the \emph{easy} one.

\begin{center}
\small
\renewcommand{\arraystretch}{1.2}
\begin{tabularx}{\linewidth}{@{}>{\raggedright\arraybackslash}p{2.6cm} >{\raggedright\arraybackslash}X >{\raggedright\arraybackslash}X@{}}
\toprule
\textbf{Dimension} & \textbf{Easy case} & \textbf{Hard case} \\
\midrule
accessible vs.\ inaccessible & the agent can obtain \emph{complete, accurate, up-to-date}
  information about the environment's state & most moderately complex environments
  (the physical world, the Internet) are inaccessible \\
deterministic vs.\ non-deterministic & any action has a \emph{single guaranteed effect}: no
  uncertainty about the resulting state & the physical world is, for all intents and
  purposes, non-deterministic \\
episodic vs.\ non-episodic & performance depends on discrete \emph{episodes} with no link
  between them: the agent decides based only on the current episode & the agent must reason
  about how current actions affect future episodes \\
static vs.\ dynamic & remains unchanged except by the agent's own actions & other processes
  operate on it and change it beyond the agent's control; the physical world is highly
  dynamic \\
discrete vs.\ continuous & a fixed, finite number of actions and percepts & the real
  world is continuous \\
\bottomrule
\end{tabularx}
\captionof{table}{Kinds of environments (slides 6--10). The more an environment is on the
left, the simpler it is to build agents for it.}
\end{center}

\begin{nbexercise}{E.1}{Practical exercise, step 1: think about the kind of environment (slide 10)}
Classify these environments along the five dimensions:
(a) a chess program (with a clock); (b) an e-mail spam filter; (c) a self-driving taxi;
(d) a robot vacuum cleaner; (e) an online poker bot.
\end{nbexercise}
\begin{nbsolution}
\small
\begin{tabular}{@{}lccccc@{}}
\toprule
 & \textbf{access.} & \textbf{determ.} & \textbf{episodic} & \textbf{static} & \textbf{discrete} \\
\midrule
(a) chess with clock & yes & yes & no & semi$^\ast$ & yes \\
(b) spam filter & no & no & yes & no & yes \\
(c) self-driving taxi & no & no & no & no & no \\
(d) robot vacuum & no & no & no & no & no \\
(e) poker bot & no & no & no$^{\ast\ast}$ & no & yes \\
\bottomrule
\end{tabular}

\normalsize
Explanations: chess is fully observable and every move has a certain effect; the board
does not change while you think, but the clock does ($^\ast$semi-dynamic). A spam filter
classifies each e-mail independently (episodic), but it cannot see the sender's intentions
(inaccessible). The taxi is the hardest case on every dimension, like the vacuum robot (the
cat moves, the floor is continuous). Poker hides the opponents' cards (inaccessible), the
deck is shuffled (non-deterministic), and ($^{\ast\ast}$) hands are episodes, but the
opponents' model of you carries over from hand to hand. \textbf{Lesson}: almost every
real application is on the ``hard'' side, which is why architectures matter.
\end{nbsolution}

\subsection{What is an agent architecture?}
\label{sec:architectures-overview}
\slidesrc{Agent architectures: 11--14}

\begin{nbdef}{Agent architecture (slide 11)}
An \textbf{architecture} proposes a particular methodology for building an autonomous agent:
\begin{itemize}[nosep]
  \item how the construction of the agent can be decomposed into the construction of a set
        of \textbf{component modules};
  \item how these modules should be made to \textbf{interact}.
\end{itemize}
These two aspects define how the \emph{sensor data} and the \emph{current internal state}
of the agent determine the \emph{actions} (effector outputs) and the \emph{future internal
state} of the agent.
\end{nbdef}

\begin{center}
\begin{tikzpicture}[font=\small, >=Stealth,
  fn/.style={draw, thick, circle, minimum size=0.9cm, fill=nbblue!10},
  st/.style={draw, thick, minimum width=0.8cm, minimum height=0.9cm, fill=nborange!12}]
  \draw[thick, nbgray, rounded corners] (0.6,-1.4) rectangle (7.4,1.5);
  \node[font=\scriptsize, nbgray] at (6.9,-1.2) {agent};
  \node (in) at (-0.8,0) {perception $p$};
  \node[fn] (f) at (2,0) {$f$};
  \node[st] (s) at (4,0) {$s$};
  \node[fn] (g) at (6,0) {$g$};
  \node (out) at (8.8,0) {action $a$};
  \draw[->, thick] (in) -- (f);
  \draw[->, thick] (f) -- (s);
  \draw[->, thick] (s) -- (g);
  \draw[->, thick] (g) -- (out);
  \draw[->, thick] (1.0,0) |- (6,1.05) -- (g.north);
  \draw[->, thick] (4.7,0) |- (2,-0.95) -- (f.south);
  \fill (1.0,0) circle (1.6pt); \fill (4.7,0) circle (1.6pt);
\end{tikzpicture}
\captionof{figure}{From perception to action (slide 12): $f$ = state update function, $s$ =
internal state, $g$ = output function. The new state depends on the old state and the
percept; the action depends on the state and the percept.}
\end{center}

In formulas, at each step $t$:
\begin{align*}
s_t &= f(s_{t-1},\, p_t) && \reason{update the internal state with what was perceived} \\
a_t &= g(s_t,\, p_t) && \reason{choose the action}
\end{align*}

\begin{nbexample}{a smart thermostat as $f$, $s$, $g$}
\notinslides
Percept $p_t$: the measured temperature. State $s_t$: the average of the last three
temperatures (to ignore a sensor glitch). Output: switch the heater on if the average is
below 20\,$^\circ$C.
\begin{center}
\small
\begin{tabular}{@{}c c c c l@{}}
\toprule
$t$ & $p_t$ & last three & $s_t$ (average) & $a_t = g(s_t, p_t)$ \\
\midrule
1 & 19 & 21, 21, 19 & 20.33 & off \\
2 & 18 & 21, 19, 18 & 19.33 & on \\
3 & 25 (glitch) & 19, 18, 25 & 20.67 & off \\
\bottomrule
\end{tabular}
\end{center}
(Initial readings 21, 21.) The internal state is what lets an agent do better than reacting
to the last percept only; how much state, and how it is used, is exactly what distinguishes
the architectures below.
\end{nbexample}

\begin{nbfact}{Main kinds of agent architectures (slide 13)}
\begin{itemize}[nosep]
  \item \textbf{Reactive} architectures: focused on fast reactions/responses to changes
        detected in the environment.
  \item \textbf{Deliberative} (symbolic) architectures: focused on long-term planning of
        actions, centred on a set of basic goals.
  \item \textbf{Hybrid} architectures: combining a reactive side and a deliberative side.
\end{itemize}
\end{nbfact}

\begin{nbexample}{reactive vs.\ deliberative: a robot that has to reach a point (slide 14)}
\begin{center}
\small
\begin{tabularx}{\linewidth}{@{}X X@{}}
\toprule
\textbf{Reactive robot} & \textbf{Deliberative robot} \\
\midrule
a sensor in the front of the robot & an explicit representation of the environment (a map) \\
change movement right/left when the sensor detects an obstacle & a planning procedure finds
  the minimal route between the current position and the destination \\
minimal computation, based on current location and destination point & high computational
  cost; possible dynamic re-planning needed \\
\bottomrule
\end{tabularx}
\end{center}
\begin{center}
\begin{tikzpicture}[scale=0.62, font=\scriptsize, >=Stealth]
  \begin{scope}
    \draw[nbgray] (0,0) rectangle (6,4);
    \fill[nbgray!50] (2,1) rectangle (3,4); \fill[nbgray!50] (4,0) rectangle (4.6,2.6);
    \node[circle, fill=nbblue, inner sep=2pt, label=below:start] (A) at (0.6,0.6) {};
    \node[star, fill=nbgreen, inner sep=1.5pt, label=above:goal] (B) at (5.4,3.4) {};
    \draw[->, thick, nborange] (A) -- (1.8,0.6) -- (1.8,0.4) -- (2.2,0.5) -- (3.4,0.4) -- (3.6,1.4) -- (3.8,2.9) -- (4.4,2.8) -- (4.9,3.2) -- (B);
    \node at (3,-0.6) {reactive: bumps and turns};
  \end{scope}
  \begin{scope}[xshift=8cm]
    \draw[nbgray] (0,0) rectangle (6,4);
    \fill[nbgray!50] (2,1) rectangle (3,4); \fill[nbgray!50] (4,0) rectangle (4.6,2.6);
    \node[circle, fill=nbblue, inner sep=2pt, label=below:start] (A) at (0.6,0.6) {};
    \node[star, fill=nbgreen, inner sep=1.5pt, label=above:goal] (B) at (5.4,3.4) {};
    \draw[->, thick, nbpurple] (A) -- (3.5,0.5) -- (3.6,3.1) -- (B);
    \node at (3,-0.6) {deliberative: planned shortest path};
  \end{scope}
\end{tikzpicture}
\captionof{figure}{Same task, two architectures. The reactive robot computes almost nothing
but wanders; the deliberative one finds the best route but must have a correct map and
re-plan if a wall moves.}
\end{center}
\end{nbexample}

\subsection{Reactive architectures}
\label{sec:reactive-architectures}
\slidesrc{Agent architectures: 15--36}

\subsubsection{Basic ideas}
\slidesrc{Agent architectures: 15--23}

\begin{nbdef}{Reactive agent (slide 15)}
Reactive agents have \emph{at most a very simple internal representation of the world},
but provide a \textbf{tight coupling of perception and action}. They follow the
\textbf{behaviour-based paradigm}: intelligence is a product of the \textbf{interaction}
between an agent and its environment.
\end{nbdef}

\paragraph{The ant colony (slides 16--18).} A single ant has very little intelligence,
computing power or reasoning abilities; the union of many ants and their interaction forms
a highly complex, structured and efficient system. The slides add pictures of other
\emph{swarms}: bees, flocks of starlings, schools of fish and rays, migrating birds, and a
\emph{Swarm Intelligence} magazine cover with insect-like robots.

\begin{nbwarning}[Prof.\ Mascardi's remark on slide 16]
``As stressed during the class, an ant colony is an example of \textbf{multiagent system},
not of a single agent architecture: I disagree with the contents of this slide.''
The ant colony shows \emph{emergent} behaviour of many simple agents (the topic of
Section~\ref{sec:abms}). An \emph{architecture} describes how \emph{one} agent is built inside.
Keep the two levels separate in the exam.
\end{nbwarning}

\begin{nbfact}{Main characteristics of reactive architectures (slides 19--21)}
\begin{enumerate}
  \item \textbf{Emergent functionality}: simple agents, simple interaction; complex
        behaviour patterns appear as a result of the dynamic interactions; the global
        behaviour is \emph{not specified a priori} (e.g.\ the movement of robots depends on
        the obstacles they meet).
  \item \textbf{Task decomposition}: agents are composed of autonomous modules, each
        managing a given task (sensing, control, computation); minimal, low-level
        communication between modules; \emph{no global world model}; \emph{no
        planning/controller agent}.
  \item \textbf{Raw data}: basic data from sensors, no complex symbolic management of data
        as in classical AI; refusal of the \textbf{physical symbol system hypothesis}, the
        basic pillar of symbolic AI: ``intelligent behaviour can only be obtained in
        symbol-processing systems''.
\end{enumerate}
\end{nbfact}

\begin{nbdef}{Reactive behaviour (slide 22)}
Each behaviour continually maps perceptual input to action output. Reactive behaviour is a
set of \textbf{action rules}
\[ \mathit{action} : S \to A \]
where $S$ is the set of states of the environment and $A$ the set of primitive actions the
agent can perform. Example (a thermostat):
\[
\mathit{action}(s) =
\begin{cases}
\text{heater off} & \text{if the temperature is OK in state } s \\
\text{heater on} & \text{otherwise}
\end{cases}
\]
\end{nbdef}

Compared with Section~\ref{sec:architectures-overview}, a purely reactive agent has
\emph{no internal state}: $a_t = g(p_t)$, the action depends only on the current percept.

\begin{center}
\begin{tikzpicture}[font=\small, >=Stealth,
  r/.style={draw, rounded corners, minimum width=1.4cm, minimum height=0.55cm, fill=#1}]
  \draw[thick, rounded corners] (0,0) rectangle (8,3.6);
  \node[font=\small\bfseries] at (4,3.3) {Agent: stimulus--response behaviours};
  \node[draw, thick, fill=nbgreen!15, minimum width=0.7cm, minimum height=3.2cm, align=center] (S) at (-0.8,1.8) {\rotatebox{90}{Sensors}};
  \node[draw, thick, fill=nbred!15, minimum width=0.7cm, minimum height=3.2cm] (E) at (8.8,1.8) {\rotatebox{90}{Effectors}};
  \foreach \y/\i in {2.6/1, 1.8/2, 0.6/n}{
    \node[r=nbblue!10] (s\i) at (2.2,\y) {State$_\i$};
    \node[r=nborange!12] (a\i) at (5.8,\y) {Action$_\i$};
    \draw[->, thick] (s\i) -- (a\i);
  }
  \node at (2.2,1.2) {$\vdots$}; \node at (5.8,1.2) {$\vdots$};
  \draw[->, thick] (S.east) -- (0,1.8);
  \draw[->, thick] (8,1.8) -- (E.west);
\end{tikzpicture}
\captionof{figure}{Basic schema of a reactive architecture (slide 23): a table of
situation$\to$action rules between sensors and effectors, no world model, no planner.}
\end{center}

\subsubsection{Rodney Brooks and the subsumption architecture}
\slidesrc{Agent architectures: 24--30}

Rodney Brooks (director of the MIT Computer Science and Artificial Intelligence Lab,
1997--2007, and co-founder of iRobot, the maker of the Roomba vacuum cleaner) is the main
critic of symbolic AI from the reactive side.

\begin{nbfact}{Brooks' three theses (slide 25)}
\begin{enumerate}[nosep]
  \item Intelligent behaviour can be generated \textbf{without explicit representations} of
        the kind that symbolic AI proposes.
  \item Intelligent behaviour can be generated \textbf{without explicit abstract reasoning}
        of the kind that symbolic AI proposes (reduced computation on sensor-like data).
  \item Intelligence is an \textbf{emergent property} of certain complex systems.
\end{enumerate}
\end{nbfact}

\begin{nbfact}{Brooks' key ideas (slides 26--27)}
\begin{itemize}[nosep]
  \item \textbf{Situatedness}: ``real'' intelligence is situated in the world. ``The world
        is its best model'': it is always up-to-date, whereas a model is an abstraction, a
        simplification that considers some characteristics and disregards others.
  \item \textbf{Embodiment}: ``real'' intelligence requires a physical body, and cannot be
        found in disembodied systems such as theorem provers or expert systems (physical
        robots).
  \item \textbf{Intelligence and emergence}: intelligent behaviour arises from the agent's
        interaction with its environment; intelligence is ``in the eye of the beholder'',
        not an innate, isolated property.
\end{itemize}
\end{nbfact}

\begin{nbdef}{Subsumption architecture (slides 28--29)}
A \textbf{subsumption architecture} is a hierarchy of task-accomplishing
\textbf{behaviours}. Each behaviour is a rather simple rule-like structure. Behaviours
\emph{compete} to exercise control over the agent, since several may be applicable at the
same time. Behaviours are organised in \textbf{layers}:
\begin{itemize}[nosep]
  \item lower layers represent more primitive kinds of behaviour (e.g.\ avoiding obstacles);
  \item higher layers represent more complex behaviours (e.g.\ identifying an object);
  \item lower layers have \textbf{precedence} over layers further up the hierarchy.
\end{itemize}
The resulting systems are, in terms of computation, extremely simple; yet some of the
robots do tasks that would be impressive if accomplished by symbolic AI systems.
\end{nbdef}

\begin{center}
\begin{tikzpicture}[font=\small, >=Stealth,
  l/.style={draw, thick, rounded corners, minimum width=5.4cm, minimum height=0.6cm, fill=#1}]
  \node[l=nbpurple!10] (l3) at (4,3.0) {level 3: plan changes to the world};
  \node[l=nbpurple!10] (l2) at (4,2.1) {level 2: explore};
  \node[l=nbpurple!10] (l1) at (4,1.2) {level 1: wander};
  \node[l=nbblue!12] (l0) at (4,0.3) {level 0: avoid obstacles};
  \node (sen) at (-0.7,0.3) {Sensors};
  \node (act) at (8.9,0.3) {Actuators};
  \draw[->, thick] (sen) -- (l0);
  \draw[->, thick] (l0) -- (act);
  \foreach \l in {l1,l2,l3}{ \draw[->, thick, nbgray] (0.7,0.3) |- (\l.west); \draw[->, thick, nbgray] (\l.east) -| (7.3,0.3); }
  \node[font=\scriptsize, align=center, nborange] at (9.0,2.1) {higher levels\\subsume (suppress)\\lower outputs};
\end{tikzpicture}
\captionof{figure}{Brooks' layered control system (``However\ldots'', slide 30, from Brooks'
original paper). Each level is a complete behaviour working directly from sensors to
actuators; new levels are added on top of working ones.}
\end{center}

\begin{nbwarning}[``However\ldots'' (slide 30): who has precedence?]
Slide 29 says that \emph{lower} layers have precedence. Slide 30 (``However\ldots'') shows the
figure of Brooks' original paper, whose caption says that control is layered ``with
\emph{higher} level layers subsuming the roles of lower level layers when they wish to take
control''. Both views exist:
\begin{itemize}[nosep]
  \item in Brooks' original robots, a higher layer can \emph{suppress} the output of a lower
        one (that is what ``subsume'' means): e.g.\ ``explore'' overrides ``wander'';
  \item in the rule-based presentations (Wooldridge's textbook, Steels' explorer below),
        the lowest-level behaviour, such as obstacle avoidance, has the \emph{highest
        priority}, because safety comes first.
\end{itemize}
In both, each layer is a complete behaviour and the composition of layers produces the
overall behaviour without a central planner.
\end{nbwarning}

\subsubsection{Steels' Mars explorer}
\slidesrc{Agent architectures: 31--33}

Luc Steels' Mars explorer system, using the subsumption architecture, achieves near-optimal
cooperative performance in a simulated ``rock gathering on Mars'' domain: \emph{explore a
distant planet and collect samples of a precious rock. The location of the samples is not
known in advance, but they tend to be clustered.} The mother-ship emits a radio signal that
gets weaker with distance, so ``travel up the gradient'' (towards a stronger signal) leads
back to the base.

\begin{nbfact}{The rules of an individual explorer (slides 32--33), in priority order}
\begin{center}
\begin{tabular}{@{}c l@{}}
(1) & \textbf{if} detect an obstacle \textbf{then} change direction \\
(2) & \textbf{if} carrying samples \textbf{and} at the base \textbf{then} drop samples \\
(3) & \textbf{if} carrying samples \textbf{and not} at the base \textbf{then} travel up gradient \\
(4) & \textbf{if} detect a sample \textbf{then} pick sample up \\
(5) & \textbf{if} true \textbf{then} move randomly
\end{tabular}
\end{center}
Obstacle avoidance (1) is the lowest-level behaviour, hence the one with the highest
priority. Rule (5) always applies, so it fires only when nothing better is applicable.
At each step the agent fires the \emph{first} rule (smallest number) whose condition holds.
\end{nbfact}

\begin{nbexample}{a trace of one explorer}
\notinslides
\small
\begin{tabularx}{\linewidth}{@{}c X c l@{}}
\toprule
\textbf{Step} & \textbf{Situation} & \textbf{Applicable rules} & \textbf{Fired} \\
\midrule
1 & empty terrain, carrying nothing & (5) & (5) move randomly \\
2 & a rock sample in front & (4), (5) & (4) pick up \\
3 & carrying a sample, far from base & (3), (5) & (3) up gradient \\
4 & carrying, a crater ahead & (1), (3), (5) & (1) change direction \\
5 & carrying, path clear & (3), (5) & (3) up gradient \\
6 & carrying, at the base & (2), (5) & (2) drop samples \\
7 & carrying nothing, at the base & (5) & (5) move randomly \\
\bottomrule
\end{tabularx}
The behaviour ``collect rocks and bring them home'' was never written anywhere: it emerges
from five local rules and the environment. (In the full system, \emph{cooperative} rules make
agents drop radioactive ``crumbs'' on the way back, so that other agents find the clusters;
this is how near-optimal cooperation emerges.)
\end{nbexample}

\subsubsection{Advantages and limitations}
\slidesrc{Agent architectures: 34--36}

\begin{center}
\small
\begin{tabularx}{\linewidth}{@{}X X@{}}
\toprule
\textbf{Advantages (slide 34)} & \textbf{Limitations (slides 35--36)} \\
\midrule
\textbf{simplicity} of individual agents &
  agents without environment models must have sufficient information available from the
  \emph{local} environment \\
\textbf{flexibility, adaptability}: ideal in very dynamic and unpredictable environments &
  decisions based on the local environment: how to take \emph{non-local} information into
  account? ``short-term'' view \\
\textbf{computational tractability}: no complex planning/reasoning procedures, no
  continuous model update &
  no long-term planning capabilities; limited applicability (games, simulations, basic
  insect-like robots) \\
\textbf{robustness} against failure: no central planning component (e.g.\ ant colony) &
  difficult to make reactive agents that \emph{learn} (dynamic evolution of rules?) \\
\textbf{elegance} &
  behaviour emerges from interactions plus environment: hard to \emph{engineer} specific
  agents (no principled methodology); hard with many behaviours (dynamics too complex to
  understand) \\
\bottomrule
\end{tabularx}
\end{center}

\begin{nbwarning}[Prof.\ Mascardi's remarks on slide 34]
\begin{itemize}[nosep]
  \item On \emph{robustness}: ``I disagree: an ant colony is a MAS, not a single agent: a
        reactive architecture is as robust (or not robust) as any other modular software
        system.''
  \item On \emph{elegance}: ``I disagree: elegance is not a scientific, engineering
        property.''
\end{itemize}
\end{nbwarning}

\begin{nbwhiteboard}[label=wb:reactive-pros-cons]{Reactive architectures: pros and cons}{written during the agent architectures lecture}
\begin{center}
\small
\begin{tabularx}{\linewidth}{@{}>{\raggedright\arraybackslash}X|>{\raggedright\arraybackslash}X@{}}
\multicolumn{1}{c|}{\textbf{\textcolor{nbgreen!70!black}{PROS}}} & \multicolumn{1}{c}{\textbf{\textcolor{nbred}{CONS}}} \\
\hline
\rule{0pt}{2.6ex}Assuming a \textbf{dynamic, unpredictable, partially observable}
  environment, reactive architectures are suitable to \textbf{cope with such changes}. &
  Reaching \textbf{long-term objectives} with \textbf{no representation of the world} and
  \textbf{no planning/reasoning} seems harder. \\[4pt]
\textbf{Simple behaviour} $\Rightarrow$ \textbf{fast computation and reaction} (safer). &
  Lack of a \textbf{consolidated and generalised engineering method}. \\[4pt]
Natural fit for \textbf{ABMS} (agent-based modelling and simulation). & \\[4pt]
\fbox{Simpler / shorter code (?)} \ \ (written in the middle, with a question mark) & \\
\end{tabularx}
\end{center}
\wbphoto{sections/agent-architectures/reactive-architectures/assets/whiteboard-8-1.jpg}
How to read it: the professor's own summary, more cautious than slide 34.
\begin{itemize}[nosep]
  \item The first pro is about the \emph{environment}: it uses the vocabulary of
        Section~\ref{sec:agents-and-environments} (dynamic, partially observable = inaccessible).
  \item ``Safer'' because a fast reaction (braking, avoiding) can matter more than an optimal
        but late plan.
  \item ABMS is listed among the pros: reactive agents are the typical building block of
        agent-based simulations (Section~\ref{sec:abms}).
  \item ``Simpler/shorter code'' is boxed with a \emph{question mark}: individual rules are
        simple, but a large set of interacting rules is not necessarily simple to write or
        to understand (compare the second limitation of slide 36).
  \item The cons match slides 35--36: no long-term view, no principled engineering method.
\end{itemize}
\end{nbwhiteboard}

\subsection{Deliberative and hybrid architectures}
\label{sec:deliberative-hybrid}
\slidesrc{Agent architectures: 37--41}

\subsubsection{Deliberative architectures}
\slidesrc{Agent architectures: 37--39}

\begin{nbdef}{Deliberative agent architecture (slide 37)}
\begin{itemize}[nosep]
  \item an \textbf{explicit symbolic model of the world};
  \item decisions are made via \textbf{logical reasoning}, based on pattern matching and
        symbolic manipulation;
  \item the \textbf{sense--plan--act} problem-solving paradigm of classical AI planning
        systems.
\end{itemize}
\end{nbdef}

\begin{center}
\begin{tikzpicture}[font=\small, >=Stealth,
  m/.style={draw, thick, rounded corners, minimum width=2.0cm, minimum height=0.9cm, fill=#1, align=center}]
  \draw[thick, rounded corners] (0,0) rectangle (8,2.2);
  \node[font=\small\bfseries] at (4,1.95) {Agent};
  \node[draw, thick, fill=nbgreen!15, minimum width=0.7cm, minimum height=2.0cm] (S) at (-0.8,1.1) {\rotatebox{90}{Sensors}};
  \node[draw, thick, fill=nbred!15, minimum width=0.7cm, minimum height=2.0cm] (E) at (8.8,1.1) {\rotatebox{90}{Effectors}};
  \node[m=nbblue!10] (W) at (1.5,0.95) {World\\model};
  \node[m=nbpurple!10] (P) at (4,0.95) {Planner};
  \node[m=nborange!12] (X) at (6.5,0.95) {Plan\\executor};
  \draw[->, thick] (S.east) -- (W);
  \draw[->, thick] (W) -- (P);
  \draw[->, thick] (P) -- (X);
  \draw[->, thick] (X) -- (E.west);
\end{tikzpicture}
\captionof{figure}{Basic deliberative architecture (slide 38): sense (update the world model),
plan, act. Shakey with its STRIPS planner (Section~\ref{sec:ai-history}) is the classic
example.}
\end{center}

\begin{nbfact}{Problems of the deliberative approach (slide 39)}
\begin{itemize}[nosep]
  \item \textbf{Dynamic world}: the symbolic world model must be updated, and the world
        changes \emph{while} planning is being done.
  \item \textbf{Representation language}: it must be expressive enough to be useful in any
        domain, but limited enough to be computationally tractable (recall FOL's
        undecidability, Section~\ref{sec:decidability}).
  \item \textbf{Classical planning} looks for complete, optimal solutions: high
        computational cost, whereas sometimes a sub-optimal, low-cost, fast reaction is
        effective.
\end{itemize}
\end{nbfact}

\begin{nbexample}{when the plan is already wrong}
\notinslides
A delivery robot plans the shortest route through a hospital corridor (2 seconds of
planning). Meanwhile a bed is pushed into the corridor. The plan was optimal for a world that
no longer exists. A purely deliberative robot either drives into the bed or must stop and
re-plan every time anything moves. This is the motivation for hybrid architectures.
\end{nbexample}

\subsubsection{Hybrid architectures}
\slidesrc{Agent architectures: 40--41}

Many researchers have argued that neither a completely deliberative nor a completely
reactive approach is suitable for building agents. They suggest \textbf{hybrid systems},
which marry the classical and the alternative approach. The obvious way is to build an
agent out of two (or more) subsystems:
\begin{itemize}[nosep]
  \item a \textbf{deliberative} one, containing a symbolic world model, which develops plans
        and makes decisions in the way proposed by symbolic AI;
  \item a \textbf{reactive} one, capable of reacting quickly to events without complex
        reasoning.
\end{itemize}

\begin{center}
\begin{tikzpicture}[font=\small, >=Stealth,
  m/.style={draw, thick, rounded corners, minimum width=1.9cm, minimum height=0.8cm, fill=#1, align=center},
  r/.style={draw, rounded corners, minimum width=1.3cm, minimum height=0.5cm, fill=#1}]
  \draw[thick, rounded corners] (0,0) rectangle (8,4.6);
  \node[font=\small\bfseries] at (4,4.35) {Agent};
  \node[draw, thick, fill=nbgreen!15, minimum width=0.7cm, minimum height=4.4cm] (S) at (-0.8,2.3) {\rotatebox{90}{Sensors}};
  \node[draw, thick, fill=nbred!15, minimum width=0.7cm, minimum height=4.4cm] (E) at (8.8,2.3) {\rotatebox{90}{Effectors}};
  % reactive
  \draw[thick, dashed, nborange, rounded corners] (0.3,2.45) rectangle (7.7,4.05);
  \node[font=\scriptsize, nborange] at (4,3.85) {reactive component};
  \node[r=nbblue!10] (s1) at (2,3.3) {State$_1$}; \node[r=nborange!12] (a1) at (6,3.3) {Action$_1$};
  \node[r=nbblue!10] (s2) at (2,2.8) {State$_n$}; \node[r=nborange!12] (a2) at (6,2.8) {Action$_n$};
  \draw[->] (s1) -- (a1); \draw[->] (s2) -- (a2);
  % deliberative
  \draw[thick, dashed, nbpurple, rounded corners] (0.3,0.25) rectangle (7.7,1.9);
  \node[font=\scriptsize, nbpurple] at (4,0.45) {deliberative component};
  \node[m=nbblue!10] (W) at (1.6,1.15) {World\\model};
  \node[m=nbpurple!10] (P) at (4,1.15) {Planner};
  \node[m=nborange!12] (X) at (6.4,1.15) {Plan\\executor};
  \draw[->, thick] (W) -- (P); \draw[->, thick] (P) -- (X);
  \draw[->, thick] (S.east |- 3.05,3.05) -- (0.3,3.05);
  \draw[->, thick] (S.east |- 1.15,1.15) -- (W);
  \draw[->, thick] (7.7,3.05) -- (E.west |- 3.05,3.05);
  \draw[->, thick] (X) -- (E.west |- 1.15,1.15);
  \draw[->, thick, nbgray] (3.0,2.45) -- node[left, font=\scriptsize]{observations} (3.0,1.9);
  \draw[<-, thick, nbgray] (5.0,2.45) -- node[right, font=\scriptsize]{modifications} (5.0,1.9);
\end{tikzpicture}
\captionof{figure}{Hybrid agent architecture (slide 41). The reactive layer handles urgent
situations immediately; the deliberative layer plans; the two exchange observations and
modifications (e.g.\ the planner can install new reactive rules, the reactive layer reports
what it did).}
\end{center}

Hybrid architectures are usually \textbf{layered}: well-known examples are
\textbf{TouringMachines} (Ferguson, 1992) and \textbf{InteRRaP} (M\"uller and Pischel,
1994), see Section~\ref{sec:agents-history}; TouringMachines reappear in the EMAS keynote
(Section~\ref{sec:engineering-mas}).

\begin{nbwhy}[Hybrid architectures in a self-driving car]
\notinslides
The emergency-braking module is reactive: ``obstacle closer than 5\,m $\Rightarrow$ brake'',
in milliseconds, no reasoning. The route planner is deliberative: it uses a map and traffic
information and re-plans when a road is closed. A middle layer turns the plan into
manoeuvres. Without the reactive layer the car is unsafe; without the deliberative layer it
never gets anywhere useful.
\end{nbwhy}


\begin{nbsummary}[Agent architectures]
\begin{itemize}[nosep]
  \item Environments can be: accessible/inaccessible, deterministic/non-deterministic,
        episodic/non-episodic, static/dynamic, discrete/continuous. The real world is the
        hardest case on every axis.
  \item An \textbf{architecture} = decomposition of the agent into modules + how they
        interact; it determines how percepts and internal state produce actions and the next
        state: $s' = f(s, p)$, $a = g(s', p)$.
  \item \textbf{Reactive}: rules $S \to A$, no world model, tight perception--action
        coupling, emergent behaviour. Brooks: intelligence without representation and without
        abstract reasoning; situatedness, embodiment, emergence; \textbf{subsumption}
        architecture (layers of behaviours competing for control).
  \item Steels' Mars explorer: 5 prioritised rules (avoid obstacles, drop samples at base,
        go home up the gradient, pick up samples, move randomly).
  \item Reactive \textbf{pros}: simple, fast, tractable, good in dynamic unpredictable
        environments. \textbf{Cons}: short-term view, no long-term planning, hard to
        learn, no principled engineering method (Whiteboard~\ref{wb:reactive-pros-cons}).
  \item \textbf{Deliberative}: explicit symbolic world model + planner + plan executor
        (sense--plan--act); problems: world changes while planning, expressiveness vs.\
        tractability, high cost.
  \item \textbf{Hybrid}: reactive component for quick reactions + deliberative component
        for planning (TouringMachines, InteRRaP).
  \item Prof.\ Mascardi's remarks: an ant colony is a MAS, not a single-agent
        architecture; robustness comes from modularity, not from reactivity; elegance is
        not an engineering property.
\end{itemize}
\end{nbsummary}
